\chapter{Hito 9: Construcción de la Capa Mart e Indicadores Analíticos en DuckDB}

\section{Objetivo}

El objetivo de este hito es construir la capa \texttt{mart} del Data Warehouse macrofinanciero.

La capa \texttt{mart} contiene vistas analíticas listas para ser consumidas por usuarios finales, reportes y dashboards. A diferencia de \texttt{staging}, que almacena datos crudos, y \texttt{core}, que contiene el modelo dimensional, la capa \texttt{mart} entrega indicadores directamente interpretables.

Al finalizar este hito el estudiante será capaz de:

\begin{itemize}
\item Crear vistas analíticas en DuckDB.
\item Calcular indicadores macrofinancieros.
\item Construir métricas como ROA, ROE, dolarización y spread cambiario.
\item Preparar tablas listas para Streamlit.
\item Integrar la construcción de marts al pipeline.
\end{itemize}

\section{Resultado esperado}

Al finalizar este hito deberán existir las siguientes vistas:

\begin{verbatim}
mart.vw_tipo_cambio_diario

mart.vw_indicadores_banca

mart.vw_hhi_bancario

mart.vw_resumen_macrofinanciero
\end{verbatim}

Estas vistas serán la base del dashboard.

\section{Arquitectura del hito}

El flujo será:

\begin{verbatim}
core.fact_tipo_cambio
core.fact_balance_mensual
core.dim_tiempo
core.dim_moneda
core.dim_entidad_financiera
|
V
mart
|
V
Streamlit
\end{verbatim}

\section{Paso 1: Crear script de marts}

Crear:

\begin{verbatim}
sql/07_marts.sql
\end{verbatim}

Este archivo contendrá todas las vistas analíticas.

\section{Paso 2: Crear vista de tipo de cambio diario}

Agregar:

\begin{lstlisting}[language=SQL]
CREATE OR REPLACE VIEW mart.vw_tipo_cambio_diario AS

SELECT
t.fecha,

```
t.anio,

t.mes,

t.anio_mes,

m.moneda_cod,

m.moneda_nombre,

f.tasa_compra,

f.tasa_venta,

f.tasa_promedio,

f.spread
```

FROM core.fact_tipo_cambio f

JOIN core.dim_tiempo t
ON t.sk_tiempo = f.sk_tiempo

JOIN core.dim_moneda m
ON m.sk_moneda = f.sk_moneda;
\end{lstlisting}

\section{Explicación}

Esta vista transforma la tabla de hechos cambiaria en una tabla amigable para análisis.

En lugar de mostrar claves como:

\begin{verbatim}
sk_tiempo
sk_moneda
\end{verbatim}

muestra variables interpretables:

\begin{verbatim}
fecha
anio_mes
moneda_cod
tasa_promedio
spread
\end{verbatim}

\section{Paso 3: Crear vista de indicadores bancarios}

Agregar:

\begin{lstlisting}[language=SQL]
CREATE OR REPLACE VIEW mart.vw_indicadores_banca AS

SELECT
t.fecha,

```
t.anio,

t.mes,

t.anio_mes,

e.entidad_cod,

e.entidad_nombre,

e.tipo_entidad,

f.activos_totales,

f.cartera_creditos,

f.depositos,

f.patrimonio,

f.resultado_neto,

f.depositos_me,

f.cartera_me,

f.resultado_neto
    / NULLIF(f.activos_totales, 0)
    AS roa,

f.resultado_neto
    / NULLIF(f.patrimonio, 0)
    AS roe,

f.depositos_me
    / NULLIF(f.depositos, 0)
    AS dolarizacion_depositos,

f.cartera_me
    / NULLIF(f.cartera_creditos, 0)
    AS dolarizacion_cartera
```

FROM core.fact_balance_mensual f

JOIN core.dim_tiempo t
ON t.sk_tiempo = f.sk_tiempo

JOIN core.dim_entidad_financiera e
ON e.sk_entidad = f.sk_entidad;
\end{lstlisting}

\section{Indicadores calculados}

La vista anterior calcula:

\begin{itemize}
\item \texttt{roa}: retorno sobre activos.
\item \texttt{roe}: retorno sobre patrimonio.
\item \texttt{dolarizacion_depositos}: proporción de depósitos en moneda extranjera.
\item \texttt{dolarizacion_cartera}: proporción de cartera en moneda extranjera.
\end{itemize}

\section{Paso 4: Crear vista de concentración bancaria}

Agregar:

\begin{lstlisting}[language=SQL]
CREATE OR REPLACE VIEW mart.vw_hhi_bancario AS

WITH cuotas AS
(
SELECT
t.fecha,

```
    t.anio_mes,

    e.entidad_cod,

    e.entidad_nombre,

    f.activos_totales,

    f.activos_totales
        / NULLIF(
            SUM(f.activos_totales)
            OVER (
                PARTITION BY t.anio_mes
            ),
            0
        )
        AS cuota_activos

FROM core.fact_balance_mensual f

JOIN core.dim_tiempo t
    ON t.sk_tiempo = f.sk_tiempo

JOIN core.dim_entidad_financiera e
    ON e.sk_entidad = f.sk_entidad
```

)

SELECT
fecha,

```
anio_mes,

SUM(
    POWER(cuota_activos, 2)
) AS hhi_activos
```

FROM cuotas

GROUP BY
fecha,
anio_mes

ORDER BY
anio_mes;
\end{lstlisting}

\section{Explicación}

El índice HHI mide concentración.

Si una sola entidad domina el mercado, el HHI se acerca a 1.

Si el mercado está muy diversificado, el HHI se acerca a 0.

En este proyecto lo calculamos usando participación en activos.

\section{Paso 5: Crear vista resumen macrofinanciero}

Agregar:

\begin{lstlisting}[language=SQL]
CREATE OR REPLACE VIEW mart.vw_resumen_macrofinanciero AS

WITH fx AS
(
SELECT
anio_mes,

```
    AVG(tasa_promedio)
        AS tasa_promedio_mes,

    AVG(spread)
        AS spread_promedio_mes

FROM mart.vw_tipo_cambio_diario

WHERE moneda_cod = 'USD'

GROUP BY
    anio_mes
```

),

banca AS
(
SELECT
anio_mes,

```
    SUM(activos_totales)
        AS activos_totales_sistema,

    SUM(depositos)
        AS depositos_sistema,

    SUM(resultado_neto)
        AS resultado_neto_sistema,

    SUM(depositos_me)
        / NULLIF(SUM(depositos), 0)
        AS dolarizacion_depositos_sistema

FROM mart.vw_indicadores_banca

GROUP BY
    anio_mes
```

)

SELECT
b.anio_mes,

```
fx.tasa_promedio_mes,

fx.spread_promedio_mes,

b.activos_totales_sistema,

b.depositos_sistema,

b.resultado_neto_sistema,

b.dolarizacion_depositos_sistema
```

FROM banca b

LEFT JOIN fx
ON fx.anio_mes = b.anio_mes

ORDER BY
b.anio_mes;
\end{lstlisting}

\section{Explicación}

Esta vista integra información cambiaria y bancaria en un solo objeto analítico.

Será útil para construir la página principal del dashboard.

\section{Paso 6: Crear script Python para construir marts}

Crear:

\begin{verbatim}
etl/load_marts.py
\end{verbatim}

Contenido:

\begin{lstlisting}[language=Python]
from db import get_connection

def run_sql_file(path, conn):

```
with open(
    path,
    "r",
    encoding="utf-8"
) as file:

    sql = file.read()

conn.execute(sql)
```

def main():

```
conn = get_connection()

run_sql_file(
    "sql/07_marts.sql",
    conn
)

conn.close()

print(
    "Marts creados correctamente."
)
```

if **name** == "**main**":
main()
\end{lstlisting}

\section{Paso 7: Ejecutar construcción de marts}

Desde la terminal:

\begin{lstlisting}[language=bash]
python etl/load_marts.py
\end{lstlisting}

Salida esperada:

\begin{verbatim}
Marts creados correctamente.
\end{verbatim}

\section{Paso 8: Crear consulta de monitoreo de marts}

Crear:

\begin{verbatim}
sql/check_marts.sql
\end{verbatim}

Contenido:

\begin{lstlisting}[language=SQL]
SELECT
table_schema,

```
table_name
```

FROM information_schema.tables

WHERE table_schema = 'mart'

ORDER BY
table_name;
\end{lstlisting}

\section{Paso 9: Crear script para revisar marts}

Crear:

\begin{verbatim}
etl/check_marts.py
\end{verbatim}

Contenido:

\begin{lstlisting}[language=Python]
from db import get_connection

conn = get_connection()

with open(
"sql/check_marts.sql",
"r",
encoding="utf-8"
) as file:

```
sql = file.read()
```

result = conn.execute(sql).fetchdf()

print(result)

conn.close()
\end{lstlisting}

Ejecutar:

\begin{lstlisting}[language=bash]
python etl/check_marts.py
\end{lstlisting}

\section{Paso 10: Inspeccionar vista de tipo de cambio}

Crear:

\begin{verbatim}
sql/inspect_vw_tipo_cambio_diario.sql
\end{verbatim}

Contenido:

\begin{lstlisting}[language=SQL]
SELECT *
FROM mart.vw_tipo_cambio_diario
ORDER BY
fecha,
moneda_cod
LIMIT 10;
\end{lstlisting}

\section{Paso 11: Inspeccionar vista bancaria}

Crear:

\begin{verbatim}
sql/inspect_vw_indicadores_banca.sql
\end{verbatim}

Contenido:

\begin{lstlisting}[language=SQL]
SELECT *
FROM mart.vw_indicadores_banca
ORDER BY
fecha,
entidad_cod
LIMIT 10;
\end{lstlisting}

\section{Paso 12: Inspeccionar resumen macrofinanciero}

Crear:

\begin{verbatim}
sql/inspect_vw_resumen_macrofinanciero.sql
\end{verbatim}

Contenido:

\begin{lstlisting}[language=SQL]
SELECT *
FROM mart.vw_resumen_macrofinanciero
ORDER BY
anio_mes;
\end{lstlisting}

\section{Paso 13: Integrar al pipeline}

Actualizar:

\begin{verbatim}
etl/run_pipeline.py
\end{verbatim}

Contenido recomendado:

\begin{lstlisting}[language=Python]
import sys

import subprocess

steps = [
"etl/load_staging.py",
"etl/load_dimensions.py",
"etl/load_facts.py",
"etl/load_marts.py"
]

for step in steps:

```
print(
    f"Ejecutando {step}..."
)

subprocess.run(
    [
        sys.executable,
        step
    ],
    check=True
)
```

print(
"Pipeline ejecutado correctamente."
)
\end{lstlisting}

\section{Paso 14: Ejecutar pipeline completo}

Desde la terminal:

\begin{lstlisting}[language=bash]
python etl/run_pipeline.py
\end{lstlisting}

Salida esperada:

\begin{verbatim}
Ejecutando etl/load_staging.py...
Ejecutando etl/load_dimensions.py...
Ejecutando etl/load_facts.py...
Ejecutando etl/load_marts.py...

Pipeline ejecutado correctamente.
\end{verbatim}

\section{Buenas prácticas}

Durante este hito seguiremos estas reglas:

\begin{enumerate}
\item El dashboard deberá conectarse a la capa \texttt{mart}, no directamente a \texttt{core}.
\item Los indicadores deben calcularse de forma reproducible.
\item Las vistas deben tener nombres claros.
\item Cada vista debe responder a una pregunta de negocio.
\item Las vistas deben ser fáciles de consumir desde Streamlit.
\end{enumerate}

\section{Checklist de validación}

\begin{itemize}
\item[$\Box$] Vista \texttt{vw_tipo_cambio_diario} creada.
\item[$\Box$] Vista \texttt{vw_indicadores_banca} creada.
\item[$\Box$] Vista \texttt{vw_hhi_bancario} creada.
\item[$\Box$] Vista \texttt{vw_resumen_macrofinanciero} creada.
\item[$\Box$] Indicadores bancarios calculados.
\item[$\Box$] Resumen macrofinanciero generado.
\item[$\Box$] Pipeline actualizado.
\end{itemize}

\section{Entregable}

Al finalizar este hito el estudiante deberá disponer de:

\begin{itemize}
\item Capa \texttt{mart} construida.
\item Vistas analíticas para Streamlit.
\item Indicadores macrofinancieros.
\item Pipeline completo hasta data marts.
\item Base lista para visualización.
\end{itemize}

\section{Próximo hito}

En el Hito 10 construiremos el dashboard en Streamlit.

Aprenderemos a:

\begin{itemize}
\item Conectar Streamlit con DuckDB.
\item Leer vistas de la capa \texttt{mart}.
\item Construir KPIs.
\item Crear gráficos interactivos.
\item Diseñar una primera aplicación macrofinanciera.
\end{itemize}
